% Figure: shear and bending-moment diagrams for a simply supported beam
% under a central point load P at span L. V is piecewise constant at
% +P/2 and -P/2; M is triangular, peaking at PL/4 at midspan.
\begin{figure}[htbp]
\centering
\begin{tikzpicture}
% loaded beam sketch
\begin{scope}[yshift=4.2cm]
  \draw[engblack, very thick] (0,0) -- (8,0);
  % supports
  \fill[engblue] (0,0) -- (-0.25,-0.45) -- (0.25,-0.45) -- cycle;
  \draw[engblue] (-0.35,-0.6) -- (0.35,-0.6);
  \fill[engblue] (8,0) circle (0.12);
  \draw[engblue] (7.65,-0.45) -- (8.35,-0.45);
  % central load
  \draw[engred, very thick, ->] (4,1.1) -- (4,0.08);
  \node[engred, anchor=south, font=\small] at (4,1.1) {$P$};
  \node[anchor=north, font=\scriptsize] at (0,-0.65) {$R_A=P/2$};
  \node[anchor=north, font=\scriptsize] at (8,-0.6) {$R_B=P/2$};
  \node[anchor=south, font=\scriptsize] at (4,-0.55) {$L/2$};
  \draw[<->, enggray] (0,-0.3) -- (4,-0.3);
\end{scope}
% shear diagram
\begin{scope}[yshift=2.0cm]
  \draw[enggray, ->] (-0.3,0) -- (8.6,0) node[anchor=west, font=\scriptsize]{$x$};
  \node[anchor=east, font=\scriptsize] at (-0.3,0.6) {$V(x)$};
  \draw[engblue, very thick] (0,0.8) -- (4,0.8) -- (4,-0.8) -- (8,-0.8);
  \draw[engblue, thick] (0,0) -- (0,0.8);
  \draw[engblue, thick] (8,-0.8) -- (8,0);
  \node[engblue, anchor=south, font=\scriptsize] at (2,0.8) {$+P/2$};
  \node[engblue, anchor=north, font=\scriptsize] at (6,-0.8) {$-P/2$};
\end{scope}
% moment diagram
\begin{scope}[yshift=-0.2cm]
  \draw[enggray, ->] (-0.3,0) -- (8.6,0) node[anchor=west, font=\scriptsize]{$x$};
  \node[anchor=east, font=\scriptsize] at (-0.3,0.6) {$M(x)$};
  \draw[engred, very thick] (0,0) -- (4,1.2) -- (8,0);
  \fill[engred!12] (0,0) -- (4,1.2) -- (8,0) -- cycle;
  \node[engred, anchor=south, font=\scriptsize] at (4,1.2) {$M_{\max}=PL/4$};
\end{scope}
\end{tikzpicture}
\caption[Shear and moment diagrams, central point load]{Loaded beam,
shear diagram, and bending-moment diagram for a simply supported span
$L$ carrying a central point load $P$. The reactions are each $P/2$.
The shear is constant at $+P/2$ left of the load and $-P/2$ to its
right, jumping by $P$ where the load acts. The moment is the running
integral of the shear, $dM/dx = V$: it rises linearly to a peak of
$PL/4$ at midspan, where the shear changes sign, then falls linearly
to zero at the right support. The peak moment sets the maximum
bending stress $\sigma_{\max}=M_{\max}/S$.}
\label{fig:vol03ch08:shear-moment}
\end{figure}
